% 图 2.12 Dy3+ (4f9) 基态的洪特规则填充图：7 个自旋向上，2 个自旋向下在 m_l=3,2
\documentclass[tikz,border=5pt]{standalone}
\usetikzlibrary{arrows.meta}
\usepackage{amsmath}
\usepackage{bm}
\usepackage{fontspec}
\setmainfont{Times New Roman}
\begin{document}
\begin{tikzpicture}[line width=0.8pt,>=Stealth]
% 表格线
\draw[line width=0.9pt] (-2.2,0.3) -- (2.05,0.3);
\draw[line width=0.9pt] (0,1.05) -- (0,-5.9);
% 表头
\node[inner sep=2pt] at (0.55,0.72) {\large$\uparrow$};
\node[inner sep=2pt] at (1.55,0.72) {\large$\downarrow$};
% 行标签与电子（圆点）
\foreach \y/\lab in {0/{{$m_l=3$}},-0.9/{{$2$}},-1.8/{{$1$}},-2.7/{{$0$}},
  -3.6/{{$-1$}},-4.5/{{$-2$}},-5.4/{{$-3$}}}{
  \node[anchor=east,inner sep=2pt] at (-0.35,\y) {\lab};
}
% 自旋向上：全部 7 行
\foreach \y in {0,-0.9,...,-5.4} \fill (0.55,\y) circle (0.09);
% 自旋向下：m_l=3 与 m_l=2
\fill (1.55,0) circle (0.09);
\fill (1.55,-0.9) circle (0.09);
\end{tikzpicture}
\end{document}
